% ============================================================
%  CHAPITRE 4 — Sprint 4 : Services additionnels
%  Application : Coffee Time
% ============================================================

\chapter{Sprint 4 : Services additionnels}

% ────────────────────────────────────────────────────────────
\section{Introduction}
% ────────────────────────────────────────────────────────────

Ce chapitre présente le quatrième et dernier sprint du projet \textbf{Coffee Time},
intitulé « Services additionnels ». Après avoir déployé les fonctionnalités de base
lors des sprints précédents — gestion du menu, prise de commande, suivi des tables et
gestion des comptes — ce sprint enrichit l'application avec des services à forte valeur
ajoutée destinés à fidéliser la clientèle et à animer son séjour en salle :
un \textbf{assistant conversationnel intelligent} basé sur la technique RAG
(\textit{Retrieval-Augmented Generation}), un \textbf{programme de fidélité} à paliers,
un \textbf{fil d'actualités} et un espace \textbf{jeux de divertissement}.

La suite du chapitre s'organise en trois parties : la planification et la conception
du sprint, illustrées par le backlog, les raffinements de cas d'utilisation et les
diagrammes de séquence ; puis la revue du sprint présentant les interfaces réalisées.

% ────────────────────────────────────────────────────────────
\section{Sprint 4 : Services additionnels}
% ────────────────────────────────────────────────────────────

\subsection{Backlog du Sprint 4}

Le tableau~\ref{tab:backlog_sprint4} présente les \textit{user stories} planifiées
pour ce sprint sous l'intitulé « Services mobiles complémentaires ».

\begin{table}[H]
\centering
\caption{Backlog du Sprint 4 — Services mobiles complémentaires}
\label{tab:backlog_sprint4}
\renewcommand{\arraystretch}{1.4}
\begin{tabular}{|c|p{3.5cm}|p{6.2cm}|c|c|}
\hline
\textbf{Id} & \textbf{Fonctionnalité} & \textbf{User Story} & \textbf{Priorité} & \textbf{Complexité} \\
\hline
US18 & Jeux de divertissement &
  En tant que client, je veux accéder à une sélection de mini-jeux thématiques
  café pour m'occuper agréablement pendant mon attente en salle. &
  Moyenne & 8 \\
\hline
US20 & Programme de fidélité &
  En tant que client fidèle, je veux cumuler des points à chaque commande,
  visualiser mon palier (Bronze, Silver, Gold, Platinum) et échanger mes points
  contre des récompenses. &
  Haute & 8 \\
\hline
US21 & Fil d'actualités &
  En tant que client, je veux consulter les dernières actualités, offres et
  événements du café directement depuis l'application mobile. &
  Moyenne & 5 \\
\hline
US22 & Assistant conversationnel &
  En tant que client, je veux poser des questions sur le menu et recevoir des
  réponses précises et vérifiées basées uniquement sur le catalogue réel du café. &
  Haute & 13 \\
\hline
\end{tabular}
\end{table}

% ────────────────────────────────────────────────────────────
\subsection{Raffinement des cas d'utilisation des services additionnels}
% ────────────────────────────────────────────────────────────

\subsubsection*{A- Cas d'utilisation « Assistant conversationnel »}

La figure~\ref{fig:uc_chatbot} présente le diagramme de cas d'utilisation
relatif à l'assistant conversationnel Coffee Time.

\begin{figure}[H]
  \centering
  \includegraphics[width=0.85\textwidth]{figures/sprint4/uc_chatbot}
  \caption{Diagramme de cas d'utilisation — Assistant conversationnel}
  \label{fig:uc_chatbot}
\end{figure}

\begin{table}[H]
\centering
\caption{Description textuelle de cas d'utilisation — Assistant conversationnel}
\label{tab:uc_chatbot}
\renewcommand{\arraystretch}{1.5}
\begin{tabular}{|p{4cm}|p{9.5cm}|}
\hline
\textbf{Cas d'utilisation} & Interagir avec l'assistant conversationnel \\
\hline
\textbf{Acteur} & Client \\
\hline
\textbf{Pré-condition} & L'application mobile est lancée et le client accède à l'onglet Assistant. \\
\hline
\textbf{Post-condition} & Le client reçoit une réponse précise et vérifiée basée sur le catalogue réel du café. \\
\hline
\textbf{Description du scénario principal} &
  \textbf{1.} Le client saisit une question sur le menu ou sélectionne une suggestion rapide. \newline
  \textbf{2.} Le système encode la question en vecteur sémantique (modèle multilingual-e5-small). \newline
  \textbf{3.} Le système récupère les produits et FAQ les plus pertinents par similarité cosinus. \newline
  \textbf{4.} Le contexte structuré est transmis au modèle de langage LLaMA 3.1 8B. \newline
  \textbf{5.} La réponse générée est validée (vérification des prix) puis retournée au client. \newline
  \textbf{6.} L'interface affiche la réponse avec les fiches produits sources. \\
\hline
\end{tabular}
\end{table}

\subsubsection*{B- Cas d'utilisation « Programme de fidélité »}

La figure~\ref{fig:uc_loyalty} présente le diagramme de cas d'utilisation
du programme de fidélité à paliers.

\begin{figure}[H]
  \centering
  \includegraphics[width=0.85\textwidth]{figures/sprint4/uc_loyalty}
  \caption{Diagramme de cas d'utilisation — Programme de fidélité}
  \label{fig:uc_loyalty}
\end{figure}

\begin{table}[H]
\centering
\caption{Description textuelle de cas d'utilisation — Programme de fidélité}
\label{tab:uc_loyalty}
\renewcommand{\arraystretch}{1.5}
\begin{tabular}{|p{4cm}|p{9.5cm}|}
\hline
\textbf{Cas d'utilisation} & Consulter le programme de fidélité \\
\hline
\textbf{Acteur} & Client \\
\hline
\textbf{Pré-condition} & Le client possède un compte fidélité Coffee Time (Bronze, Silver, Gold ou Platinum). \\
\hline
\textbf{Post-condition} & Le client consulte son solde de points, son palier et peut échanger ses points contre une récompense. \\
\hline
\textbf{Description du scénario principal} &
  \textbf{1.} Le client navigue vers l'onglet « Fidélité » de l'application mobile. \newline
  \textbf{2.} Le système affiche la carte de fidélité avec le palier actuel, le solde de points disponibles et la barre de progression vers le palier suivant. \newline
  \textbf{3.} Le client consulte l'historique de ses transactions de points. \newline
  \textbf{4.} Le client accède au catalogue des récompenses et sélectionne une offre à échanger. \newline
  \textbf{5.} Le système vérifie que le solde est suffisant et procède à l'échange. \newline
  \textbf{6.} Le solde est mis à jour et une confirmation est affichée. \\
\hline
\end{tabular}
\end{table}

\subsubsection*{C- Cas d'utilisation « Jeux de divertissement »}

La figure~\ref{fig:uc_games} présente le diagramme de cas d'utilisation
de l'espace jeux de divertissement.

\begin{figure}[H]
  \centering
  \includegraphics[width=0.85\textwidth]{figures/sprint4/uc_games}
  \caption{Diagramme de cas d'utilisation — Jeux de divertissement}
  \label{fig:uc_games}
\end{figure}

\begin{table}[H]
\centering
\caption{Description textuelle de cas d'utilisation — Jeux de divertissement}
\label{tab:uc_games}
\renewcommand{\arraystretch}{1.5}
\begin{tabular}{|p{4cm}|p{9.5cm}|}
\hline
\textbf{Cas d'utilisation} & Jouer à un jeu de divertissement \\
\hline
\textbf{Acteur} & Client \\
\hline
\textbf{Pré-condition} & Le client est en salle et accède à l'onglet « Fun » de l'application. \\
\hline
\textbf{Post-condition} & Le client a participé à un mini-jeu et son score est enregistré localement. \\
\hline
\textbf{Description du scénario principal} &
  \textbf{1.} Le client ouvre l'onglet « Fun » et visualise la grille des six mini-jeux disponibles. \newline
  \textbf{2.} Le client sélectionne un jeu (Quiz Barista, Mots Cachés, Mots Croisés, Word Brew, Memory ou Coffee Quiz). \newline
  \textbf{3.} Le jeu se lance et affiche les instructions ainsi que le score courant. \newline
  \textbf{4.} Le client joue une ou plusieurs parties. \newline
  \textbf{5.} À la fin d'une partie, le système enregistre le score et l'affiche dans le récapitulatif. \newline
  \textbf{6.} Le client peut rejouer ou choisir un autre jeu. \\
\hline
\end{tabular}
\end{table}

% ────────────────────────────────────────────────────────────
\subsection{Diagrammes de séquence}
% ────────────────────────────────────────────────────────────

\subsubsection*{A- Diagramme de séquence « Interagir avec l'assistant »}

La figure~\ref{fig:seq_chatbot} illustre le flux complet d'une interaction avec
l'assistant conversationnel, depuis l'envoi du message par le client jusqu'à la
réception de la réponse validée.

\begin{figure}[H]
  \centering
  \includegraphics[width=\textwidth]{figures/sprint4/seq_chatbot}
  \caption{Diagramme de séquence — Interagir avec l'assistant}
  \label{fig:seq_chatbot}
\end{figure}

Le message du client est reçu par le backend via la route
\texttt{POST /api/chat/message}. Le \texttt{RagService} encode la question en
vecteur sémantique à l'aide du modèle \textbf{multilingual-e5-small}, puis compare
ce vecteur à ceux des produits actifs stockés dans PostgreSQL par similarité cosinus.
Les cinq produits et les deux FAQ les plus proches sont sélectionnés et injectés
comme contexte dans la requête envoyée à l'API Groq (\textbf{LLaMA 3.1 8B Instant},
\texttt{temperature=0.1}). La réponse générée est ensuite vérifiée : tout prix
mentionné est comparé aux données réelles du contexte avant d'être retourné
au client avec les fiches sources.

\subsubsection*{B- Diagramme de séquence « Consulter le programme de fidélité »}

La figure~\ref{fig:seq_loyalty} illustre le flux de consultation et d'attribution
de points lors du passage en caisse.

\begin{figure}[H]
  \centering
  \includegraphics[width=0.85\textwidth]{figures/sprint4/seq_loyalty}
  \caption{Diagramme de séquence — Consulter le programme de fidélité}
  \label{fig:seq_loyalty}
\end{figure}

Lors du passage en caisse, le système calcule les points à attribuer au taux de
\textbf{1 point par tranche de 50 TND} dépensés. Le \texttt{LoyaltyService} persiste
la transaction dans la table \texttt{loyalty\_transactions} et met à jour le solde
du compte. Le palier (Bronze, Silver à 500 pts, Gold à 1\,500 pts, Platinum à
4\,000 pts) est recalculé côté client à partir du total de points cumulés à vie,
ce qui permet un affichage instantané sans requête réseau supplémentaire.

% ────────────────────────────────────────────────────────────
\section{Revue du Sprint 4}
% ────────────────────────────────────────────────────────────

\subsection{Interfaces mobiles}

Cette section présente les interfaces réalisées au cours du Sprint 4,
regroupées par fonctionnalité.

\paragraph{Assistant conversationnel.}

\begin{figure}[H]
  \centering
  \begin{minipage}[b]{0.30\textwidth}
    \centering
    \includegraphics[width=\textwidth]{figures/sprint4/chat_welcome}
    \caption*{(a) Accueil assistant}
  \end{minipage}
  \hfill
  \begin{minipage}[b]{0.30\textwidth}
    \centering
    \includegraphics[width=\textwidth]{figures/sprint4/chat_conversation}
    \caption*{(b) Conversation}
  \end{minipage}
  \hfill
  \begin{minipage}[b]{0.30\textwidth}
    \centering
    \includegraphics[width=\textwidth]{figures/sprint4/chat_sources}
    \caption*{(c) Fiches sources}
  \end{minipage}
  \caption{Interfaces de l'assistant conversationnel Coffee Time}
  \label{fig:chat_ui}
\end{figure}

L'interface de l'assistant conversationnel reprend la charte graphique
espresso/crème de l'application. L'en-tête affiche le nom « Coffee Time »,
un indicateur de statut en ligne et un bouton de réinitialisation. Les messages
du client sont présentés dans des bulles à dégradé sombre, tandis que les
réponses de l'assistant apparaissent sur fond blanc avec un avatar circulaire
et un indicateur de frappe animé. Une barre de suggestions rapides, mise à jour
dynamiquement selon le contexte, facilite la navigation.

\paragraph{Programme de fidélité.}

\begin{figure}[H]
  \centering
  \begin{minipage}[b]{0.30\textwidth}
    \centering
    \includegraphics[width=\textwidth]{figures/sprint4/loyalty_card}
    \caption*{(a) Carte et solde}
  \end{minipage}
  \hfill
  \begin{minipage}[b]{0.30\textwidth}
    \centering
    \includegraphics[width=\textwidth]{figures/sprint4/loyalty_rewards}
    \caption*{(b) Catalogue récompenses}
  \end{minipage}
  \hfill
  \begin{minipage}[b]{0.30\textwidth}
    \centering
    \includegraphics[width=\textwidth]{figures/sprint4/loyalty_history}
    \caption*{(c) Historique}
  \end{minipage}
  \caption{Interfaces du programme de fidélité}
  \label{fig:loyalty_ui}
\end{figure}

L'écran de fidélité affiche une carte animée indiquant le palier actuel du client
(Bronze, Silver, Gold ou Platinum) avec son code couleur distinctif, le solde de
points disponibles et une barre de progression vers le palier suivant. Un onglet
\textit{Récompenses} liste le catalogue échangeable (espresso offert à 150 pts,
réduction de 10\% à 300 pts, boisson gratuite à 500 pts, repas offert à 1\,000 pts).

\paragraph{Jeux de divertissement.}

\begin{figure}[H]
  \centering
  \begin{minipage}[b]{0.45\textwidth}
    \centering
    \includegraphics[width=\textwidth]{figures/sprint4/games_grid}
    \caption*{(a) Grille des jeux}
  \end{minipage}
  \hfill
  \begin{minipage}[b]{0.45\textwidth}
    \centering
    \includegraphics[width=\textwidth]{figures/sprint4/games_quiz}
    \caption*{(b) Quiz Barista}
  \end{minipage}
  \caption{Interfaces de l'espace jeux de divertissement}
  \label{fig:games_ui}
\end{figure}

L'espace « Fun » propose six mini-jeux thématiques café accessibles en un tap
depuis la grille de l'onglet dédié : Quiz Barista, Mots Cachés, Mots Croisés,
Word Brew, Memory et Coffee Quiz. Chaque jeu est implémenté comme une page
autonome dans \texttt{mobile/app/games/}, les scores étant persistés localement
via \texttt{AsyncStorage} et gérés centralement par le \texttt{GameContext}.

\paragraph{Fil d'actualités.}

\begin{figure}[H]
  \centering
  \begin{minipage}[b]{0.45\textwidth}
    \centering
    \includegraphics[width=\textwidth]{figures/sprint4/news_feed}
    \caption*{(a) Fil d'actualités}
  \end{minipage}
  \hfill
  \begin{minipage}[b]{0.45\textwidth}
    \centering
    \includegraphics[width=\textwidth]{figures/sprint4/news_detail}
    \caption*{(b) Détail d'une actualité}
  \end{minipage}
  \caption{Interface du fil d'actualités}
  \label{fig:news_ui}
\end{figure}

Le fil d'actualités permet aux clients de consulter les dernières offres, promotions
et événements du café directement depuis l'application. Les articles sont publiés
par l'administrateur depuis le tableau de bord et affichés en temps réel dans
un flux paginé sur l'interface mobile.

\subsection{Interfaces d'administration}

\begin{figure}[H]
  \centering
  \includegraphics[width=0.92\textwidth]{figures/sprint4/loyalty_admin}
  \caption{Gestion du programme de fidélité — Interface administrateur}
  \label{fig:loyalty_admin}
\end{figure}

Le tableau de bord administrateur intègre une page \textit{Programme de Fidélité}
permettant de rechercher un compte client par nom ou numéro de téléphone,
de consulter son historique de transactions, d'ajuster manuellement un solde de
points et de visualiser des statistiques globales (répartition par palier, points
distribués ce mois-ci).

% ────────────────────────────────────────────────────────────
\section{Conclusion}
% ────────────────────────────────────────────────────────────

Ce quatrième sprint clôt le cycle de développement de l'application
\textbf{Coffee Time} en déployant les services mobiles complémentaires qui
différencient la solution sur le marché. L'assistant conversationnel basé sur
l'architecture RAG garantit des réponses précises et vérifiées grâce au système
anti-hallucination combinant un prompt système strict, une température basse et
une validation post-génération des prix. Le programme de fidélité à quatre paliers
renforce la relation client sur le long terme. Le fil d'actualités maintient
l'engagement des utilisateurs entre leurs visites, et l'espace de six mini-jeux
thématiques enrichit agréablement leur séjour en salle.

L'ensemble de ces fonctionnalités, combinées à celles des sprints précédents,
constitue une application mobile et un tableau de bord d'administration complets,
cohérents et prêts à être déployés en conditions réelles.
